\chapter{Material e Métodos}
\label{cap:material-e-metodos}

% Descreva aqui os materiais utilizados e os métodos empregados
% na realização da pesquisa. Inclua informações sobre:
% - Área de estudo / local do experimento
% - Materiais e equipamentos
% - Delineamento experimental
% - Procedimentos analíticos
% - Análises estatísticas

% [Insira o texto de material e métodos aqui]
Este capítulo descreve os materiais e os métodos empregados na condução do estudo. A apresentação segue as recomendações da \gls{ABNT}, de modo a assegurar a padronização e a reprodutibilidade dos procedimentos.

\section{Área de Estudo}
\label{sec:area-de-estudo}

% [Descreva a área ou local onde o trabalho foi realizado]
O estudo foi conduzido em ambiente representativo do contexto investigado, selecionado de acordo com critérios previamente definidos.

\section{Procedimentos Metodológicos}
\label{sec:procedimentos-metodologicos}

% [Descreva os procedimentos utilizados]
Os procedimentos adotados são sintetizados no \autoref{alg:exemplo}, que descreve, em alto nível, as etapas executadas ao longo do estudo.

\begin{algorithm}[!htbp]
    \SetSpacedAlgorithm
    \caption{\label{alg:exemplo}Procedimento geral adotado no estudo}
    \Entrada{Conjunto de dados brutos $D$}
    \Saida{Resultados consolidados $R$}
    \Inicio{
        \Para{cada elemento $d \in D$}{
            processar $d$\;
            armazenar o resultado parcial\;
        }
        consolidar os resultados em $R$\;
        \Retorna $R$\;
    }
\end{algorithm}

% Dica: Para inserir um algoritmo, use o seguinte modelo:
%
% \begin{algorithm}[h!]
%     \SetSpacedAlgorithm
%     \caption{\label{alg:modelo}Descrição do Algoritmo}
%     \Entrada{Entrada do Algoritmo}
%     \Saida{Saída do Algoritmo}
%     \Inicio{
%         Passo 1\;
%         Passo 2\;
%     }
% \end{algorithm}

\section{Análise dos Dados}
\label{sec:analise-dos-dados}

% [Descreva como os dados foram analisados]
Os dados foram analisados por meio de procedimentos estatísticos descritivos e inferenciais. A relação entre as variáveis é expressa, de forma simplificada, pela \autoref{eq:modelo}.

\begin{equation}
    \label{eq:modelo}
    y = ax + b
\end{equation}

\noindent onde $y$ representa a variável resposta, $x$ a variável explicativa, e $a$ e $b$ os coeficientes do modelo.

Para a inferência sobre a média populacional $\mu$, adotou-se o nível de significância $\alpha$ e construiu-se o intervalo de confiança a partir do desvio-padrão $\sigma$ e do tamanho da amostra $n$, conforme a \autoref{eq:intervalo}.

\begin{equation}
    \label{eq:intervalo}
    \mu = \bar{x} \pm z_{\alpha/2}\,\frac{\sigma}{\sqrt{n}}
\end{equation}

% Dica: Para inserir uma fórmula matemática numerada e referenciável:
%
% \begin{equation}
%     \label{eq:exemplo}
%     y = ax + b
% \end{equation}
%
% Referencie no texto com \autoref{eq:exemplo} (ou \autoref{sec:area-de-estudo}).
